% =====================================================================
% Fundamentos de Inteligencia Artificial -- 1o Trabalho
% Mundo dos Blocos de Tamanho Variavel via SAT
%
% COMO USAR (Overleaf): New Project > Upload Project > envie o .zip.
%   - Cada secao esta num arquivo separado em secoes/. Cada membro edita
%     so o seu arquivo, e ninguem atrapalha o outro.
%   - \afazer{...} marca em vermelho o que falta. Apague quando terminar.
%   - \responsavel{Nome} indica quem escreve a secao.
%   - Exemplos prontos (desenho de estado, tabela de passos, grafo de
%     ordem parcial, saida de terminal): descomente a linha
%     % =====================================================================
% EXEMPLOS DE USO DOS COMANDOS  (nao faz parte da entrega)
% Copie o trecho que precisar para a sua secao.
% =====================================================================
\newpage
\section*{Exemplos de uso (apagar antes de entregar)}

\subsection*{1. Desenhar um estado}
\begin{figure}[H]
\centering
\begin{tikzpicture}
  % \bl{ponto p}{nivel l}{comprimento}{nome}
  \mesa
  \bl{0}{0}{2}{c}\bl{3}{0}{1}{a}\bl{5}{0}{1}{b}\bl{3}{1}{3}{d}
  \rot{$S_0$}
\end{tikzpicture}
\caption{Um estado.}
\end{figure}

\subsection*{2. Dois estados com uma ação entre eles}
\begin{figure}[H]
\centering
\begin{tikzpicture}
  \begin{scope}\mesa
    \bl{0}{0}{2}{c}\bl{3}{0}{1}{a}\bl{5}{0}{1}{b}\bl{3}{1}{3}{d}\rot{$S_0$}
  \end{scope}
  \draw[-{Stealth}] (3.3,0.5) -- node[above,font=\scriptsize]{move(d,c,0)} (5.1,0.5);
  \begin{scope}[xshift=5.4cm]\mesa
    \bl{0}{0}{2}{c}\bl{3}{0}{1}{a}\bl{5}{0}{1}{b}\bl{0}{1}{3}{d}\rot{$S_1$}
  \end{scope}
\end{tikzpicture}
\caption{Uma transição.}
\end{figure}

\subsection*{3. Tabela de passos (ação, pré-condições, ADD, DEL)}
\begin{table}[H]
\centering\small
\renewcommand{\arraystretch}{1.25}
\begin{tabularx}{\textwidth}{@{}p{2.3cm} X X X@{}}
\toprule
Ação & Pré-condições (verificadas) & ADD & DEL\\
\midrule
$A_1$: move(d,c,0)\newline $t=0$ &
$\clr(d)$; \dots &
$\at(d,0)$, $\on(d,c)$, \dots &
$\at(d,3)$, $\on(d,a)$, \dots\\
\bottomrule
\end{tabularx}
\caption{Execução manual passo a passo.}
\end{table}

\subsection*{4. Fórmulas}
\begin{align}
&\clr(b,t) \leftrightarrow \forall s,l\,\bigl(\cov(b,s,l,t)\rightarrow
  \esp(b,s,t)\bigr) \tag{A5b}
\end{align}
Vínculo causal no texto: $A_1\xrightarrow{\;\clr(a)\;}A_2$. Ordem: $A_1\prec A_2$.

\subsection*{5. Grafo de ordem parcial}
\begin{figure}[H]
\centering
\begin{tikzpicture}[node distance=1.0cm,
  acao/.style={draw,rounded corners,fill=cora!12,minimum width=4.4cm,
               minimum height=0.75cm,font=\small},
  lig/.style={-{Stealth},thick}]
  \node[acao,fill=black!6] (ini) {Início};
  \node[acao,below=of ini] (a1) {$A_1$: move(d,c,0)};
  \node[acao,below=of a1,fill=black!6] (fim) {Fim};
  \draw[lig] (ini) -- node[right,font=\footnotesize]{\ $\clr(d)$} (a1);
  \draw[lig] (a1) -- node[right,font=\footnotesize]{\ \dots} (fim);
\end{tikzpicture}
\caption{Grafo de ordem parcial.}
\end{figure}

\subsection*{6. Comandos e saída de terminal}
\begin{lstlisting}
$ python3 ../bw2cnf_var.py --cenario 1
Gerado: ... variaveis, ... clausulas (T=4)
\end{lstlisting}

\subsection*{7. Tabela simples}
\begin{center}\small
\begin{tabular}{@{}lll@{}}
\toprule
Coluna 1 & Coluna 2 & Coluna 3\\
\midrule
a & b & c\\
\bottomrule
\end{tabular}
\end{center}
 no fim deste arquivo.
% =====================================================================
\documentclass[11pt,a4paper]{article}
% =====================================================================
% preambulo.tex -- pacotes e comandos comuns. NAO precisa editar.
% =====================================================================
\usepackage[utf8]{inputenc}
\usepackage[T1]{fontenc}
\usepackage[brazil]{babel}
\usepackage{lmodern}
\usepackage[margin=2.3cm]{geometry}
\usepackage{amsmath,amssymb}
\usepackage{tikz}
\usetikzlibrary{arrows.meta,positioning}
\usepackage{booktabs,tabularx,array}
\usepackage{xcolor}
\usepackage{listings}
\usepackage{float}
\usepackage{caption}
\usepackage[hidelinks]{hyperref}
\setlength{\emergencystretch}{3em}

% ---------- cores dos blocos ----------
\definecolor{cora}{RGB}{140,90,255}
\definecolor{corb}{RGB}{255,240,0}
\definecolor{corc}{RGB}{0,170,255}
\definecolor{cord}{RGB}{235,30,30}

% ---------- blocos de codigo / saida de terminal ----------
\lstset{basicstyle=\ttfamily\footnotesize,frame=single,breaklines=true,
  columns=fullflexible,keepspaces=true,backgroundcolor=\color{black!3},
  rulecolor=\color{black!30}}

% ---------- marcadores para o trabalho em equipe ----------
% \afazer{texto}      -> aviso vermelho no PDF do que ainda falta escrever
% \responsavel{Nome}  -> quem escreve a secao (aparece em cinza)
% Para sumir com os dois na versao final, troque \marcadorestrue por
% \marcadoresfalse no main.tex.
\newif\ifmarcadores
\newcommand{\afazer}[1]{\ifmarcadores\par\noindent
  \textcolor{red!75!black}{\textbf{[A FAZER]} #1}\par\fi}
\newcommand{\responsavel}[1]{\ifmarcadores\par\noindent
  \textcolor{black!50}{\small\itshape Responsável: #1}\par\smallskip\fi}

% ---------- desenho de estados (TikZ) ----------
% Dentro de um tikzpicture:
%   \mesa                      desenha a mesa com os pontos 0..6
%   \bl{p}{l}{comprimento}{nome}   bloco comecando no ponto p, nivel l
%   \rot{texto}                rotulo acima do estado (ate o nivel 1)
%   \rotA{texto}               rotulo mais alto (estados com nivel 2)

% ---------- desenho de estados ----------
\newcommand{\mesa}{%
	\draw[line width=0.7pt] (0,0) -- (3,0);
	\foreach \i in {0,...,6}{\draw (\i*0.5,0) -- (\i*0.5,-0.07);
		\node[below,inner sep=1pt,font=\tiny] at (\i*0.5,-0.07) {\i};}}
% \bl{ponto p}{nivel l}{comprimento}{nome}
\newcommand{\bl}[4]{%
	\filldraw[fill=cor#4,draw=black!70] (#1*0.5,#2*0.5) rectangle ++(#3*0.5,0.5);
	\node[font=\small] at (#1*0.5+#3*0.25,#2*0.5+0.25) {\textit{#4}};}
\newcommand{\rot}[1]{\node[font=\small] at (1.5,1.35) {#1};}

\newcommand{\Szero}{\bl{0}{0}{2}{c}\bl{3}{0}{1}{a}\bl{5}{0}{1}{b}\bl{3}{1}{3}{d}}
\newcommand{\Sum}{\bl{0}{0}{2}{c}\bl{3}{0}{1}{a}\bl{5}{0}{1}{b}\bl{0}{1}{3}{d}}
\newcommand{\Sdois}{\bl{0}{0}{2}{c}\bl{0}{1}{3}{d}\bl{5}{0}{1}{b}\bl{5}{1}{1}{a}}
\newcommand{\Stres}{\bl{0}{0}{2}{c}\bl{2}{0}{3}{d}\bl{5}{0}{1}{b}\bl{5}{1}{1}{a}}
\newcommand{\Sfq}{\bl{0}{0}{2}{c}\bl{2}{0}{3}{d}\bl{5}{0}{1}{b}\bl{0}{1}{1}{a}}
\newcommand{\rotA}[1]{\node[font=\small] at (1.5,1.85) {#1};}
% Situacao 2
\newcommand{\Tzero}{\bl{0}{0}{2}{c}\bl{3}{0}{3}{d}\bl{0}{1}{1}{a}\bl{1}{1}{1}{b}}
\newcommand{\Tum}{\bl{0}{0}{2}{c}\bl{2}{0}{1}{b}\bl{3}{0}{3}{d}\bl{0}{1}{1}{a}}
\newcommand{\Tdois}{\bl{0}{0}{2}{c}\bl{2}{0}{1}{b}\bl{3}{0}{3}{d}\bl{2}{1}{1}{a}}
\newcommand{\Ttres}{\bl{2}{0}{1}{b}\bl{3}{0}{3}{d}\bl{2}{1}{1}{a}\bl{4}{1}{2}{c}}
\newcommand{\Tquatro}{\bl{2}{0}{1}{b}\bl{3}{0}{3}{d}\bl{4}{1}{2}{c}\bl{4}{2}{1}{a}}
\newcommand{\Tcinco}{\bl{3}{0}{3}{d}\bl{4}{1}{2}{c}\bl{4}{2}{1}{a}\bl{5}{2}{1}{b}}
% Situacao 3
\newcommand{\Ucinco}{\bl{0}{0}{2}{c}\bl{2}{0}{3}{d}\bl{0}{1}{1}{a}\bl{1}{1}{1}{b}}
\newcommand{\Usete}{\bl{0}{0}{2}{c}\bl{3}{0}{3}{d}\bl{0}{1}{1}{a}\bl{1}{1}{1}{b}}

% ---------- predicados (modo matematico) ----------
\newcommand{\at}{\mathit{at}}
\newcommand{\lev}{\mathit{lev}}
\newcommand{\clr}{\mathit{clr}}
\newcommand{\cov}{\mathit{cov}}
\newcommand{\livre}{\mathit{livre}}
\newcommand{\esp}{\mathit{esp}}
\newcommand{\on}{\mathit{on}}
\newcommand{\mv}{\mathit{move}}
\newcommand{\spn}{\mathit{span}}
\newcommand{\ov}{\mathit{overlap}}
\newcommand{\mesaT}{\mathsf{T}}


\marcadorestrue      % troque por \marcadoresfalse na versao final

\title{Mundo dos Blocos de Tamanho Variável:\\
Planejamento via SAT --- Situações 1, 2 e 3}
\author{Equipe XX --- Fundamentos de Inteligência Artificial (IComp/UFAM)\\[4pt]
Nome do Integrante 1 \and Nome do Integrante 2 \and Nome do Integrante 3}
\date{Outubro de 2026}

\begin{document}
\maketitle

% ---------- divisao de tarefas (apagar na versao final, se quiserem) ----------
\ifmarcadores
\begin{center}\small
\begin{tabular}{@{}lll@{}}
\toprule
Seção & Arquivo & Responsável\\
\midrule
1. Introdução ao Problema & \texttt{secoes/01\_introducao.tex} & \\
2. Descrição Formal & \texttt{secoes/02\_descricao\_formal.tex} & \\
3. Codificação CNF & \texttt{secoes/03\_codificacao\_cnf.tex} & \\
4. Exemplos dos 3 cenários & \texttt{secoes/04\_exemplos.tex} & \\
5. Mapeamento: formal $\to$ código & \texttt{secoes/05\_mapeamento.tex} & \\
6. Execução passo a passo & \texttt{secoes/06\_execucao.tex} & \\
7. Interpretação da saída do SAT & \texttt{secoes/07\_interpretacao.tex} & \\
\bottomrule
\end{tabular}
\end{center}
\fi

\tableofcontents
\newpage

% =====================================================================
\section{Introdução ao Problema}
\label{sec:intro}
% =====================================================================
\responsavel{(nome)}

\afazer{Explicar o Mundo dos Blocos clássico e o que muda com blocos de tamanhos
diferentes (espaços vazios, pontes, equilíbrio, movimentos laterais).}

\afazer{Dizer o que é o problema de planejamento e as etapas do trabalho: descrição
em LPO, plano manual, ordem parcial, codificação CNF e SAT.}

\afazer{Escopo: quais cenários foram resolvidos (Situação 1 até $S_{f?}$, Situação 2
até $S_5$, Situação 3 até $S_7$) e uma figura com o estado inicial e a meta de cada um.}

\afazer{Declarar o uso do chatbot de IA, como pede o enunciado.}

% =====================================================================
\section{Descrição Formal do Mundo dos Blocos de Tamanho Variado}
\label{sec:formal}
% =====================================================================
\responsavel{(nome)}

\subsection{Domínio}
\afazer{Blocos e comprimentos, pontos e slots, níveis, instantes, $\spn(b,p)$.}

\subsection{Predicados e axiomas}
\afazer{Tabela dos predicados ($\at$, $\lev$, $\cov$, $\livre$, $\esp$, $\clr$, $\on$) e
os axiomas em LPO (unicidade, exclusão, estabilidade etc.).}

\subsection{A ação \texorpdfstring{$\mv(b,y,p)$}{move(b,y,p)}}
\afazer{Pré-condições (P1, P2, \dots) e efeitos em estilo STRIPS (ADD / DEL).}

\subsection{Elementos novos e as ações associadas}
\afazer{Item 2 do enunciado: para cada elemento novo, o que a ação adiciona e remove.}

\subsection{Ajustes em relação ao manual}
\afazer{O que foi mudado em relação ao manual e por quê (predicado $\esp$ no lugar de
$\clr(y)$; regra da sombra).}

\subsection{Execução manual}
\label{sec:manual}
\subsubsection{Situação 1}
\responsavel{(nome)}
\afazer{Figura com os estados do plano e tabela: ação, pré-condições, ADD, DEL.
Argumento de por que o plano é mínimo.}

\subsubsection{Situação 2}
\responsavel{(nome)}
\afazer{Idem para $S_0 \to S_5$.}

\subsubsection{Situação 3}
\responsavel{(nome)}
\afazer{Idem para $S_0 \to S_7$.}

\subsection{Ordem parcial}
\label{sec:pop}
\afazer{Definição: ações, vínculos causais $A\xrightarrow{\;p\;}B$, restrições de
ordem, ameaças.}

\subsubsection{Situação 1}
\afazer{Tabela de vínculos causais, grafo, análise de ameaças, linearizações.}

\subsubsection{Situação 2}
\afazer{Idem.}

\subsubsection{Situação 3}
\afazer{Idem.}

\subsubsection{Ordem parcial no SAT}
\afazer{Como as restrições $\varphi_1\prec_P\varphi_2$ viram cláusulas e o resultado
de cada cenário.}

% =====================================================================
\section{Codificação CNF}
\label{sec:cnf}
% =====================================================================
\responsavel{(nome)}

\subsection{Variáveis proposicionais}
\afazer{Tabela: variável, índices, papel e quantidade em cada cenário.}

\subsection{Grupos de cláusulas}
\afazer{Tabela: grupo (3.1, 3.2, 3.3 (A), \dots), esquema das cláusulas e quantidade.
Os números estão no cabeçalho de cada \texttt{trab01\_blocos2SAT.cnf}.}

\subsection{Relação \texorpdfstring{$\on$}{on}: derivada, não codificada}
\afazer{Explicar que $\on$ é reconstruída pelo \texttt{interpretar.py}.}

% =====================================================================
\section{Exemplos dos 3 cenários codificados para CNF em LP}
\label{sec:exemplos}
% =====================================================================
\responsavel{(nome)}

\subsection{Situação 1}
\afazer{Estado inicial e meta como fórmulas e como cláusulas unitárias DIMACS (os ids
estão no arquivo \texttt{.map}). Amostra de cláusulas reais de alguns grupos.}

\subsection{Situação 2}
\afazer{Estado inicial e meta em fórmulas e DIMACS.}

\subsection{Situação 3}
\afazer{Estado inicial e meta em fórmulas e DIMACS.}

% =====================================================================
\section{Mapeamento: Descrição Formal \texorpdfstring{$\to$}{->} Código}
\label{sec:mapeamento}
% =====================================================================
\responsavel{(nome)}

\afazer{Tabela de duas colunas: regra formal (axioma, pré-condição, efeito) e onde
ela está no \texttt{bw2cnf\_var.py} (nome da variável, função ou bloco \texttt{\# 3.x}).}

% =====================================================================
\section{Execução Passo a Passo}
\label{sec:exec}
% =====================================================================
\responsavel{(nome)}

\afazer{Como configurar o cenário, os comandos usados e a saída de cada um: busca do
menor horizonte, geração do CNF, miniSAT e interpretador. Uma tabela com o resultado
dos 3 cenários (horizonte, variáveis, cláusulas, tempo).}

% =====================================================================
\section{Interpretação da Saída do SAT Solver}
\label{sec:interp}
% =====================================================================
\responsavel{(nome)}

\subsection{De inteiros para o plano}
\afazer{Formato do \texttt{resultadoN.txt}; como o \texttt{.map} traduz os inteiros;
filtragem dos literais positivos; ordenação das ações por $t$.}

\subsection{Comparação: plano manual \texorpdfstring{$\times$}{x} chatbot \texorpdfstring{$\times$}{x} SAT}
\afazer{Tabela comparando os planos dos 3 cenários e comentário sobre as diferenças.}

\subsection{Validação da codificação}
\afazer{Como conferimos que a CNF está correta.}

\subsection{Resumo}
\afazer{Principais resultados em 3 ou 4 itens.}


% Exemplos de uso dos comandos (descomente para ver; nao entra na entrega):
% % =====================================================================
% EXEMPLOS DE USO DOS COMANDOS  (nao faz parte da entrega)
% Copie o trecho que precisar para a sua secao.
% =====================================================================
\newpage
\section*{Exemplos de uso (apagar antes de entregar)}

\subsection*{1. Desenhar um estado}
\begin{figure}[H]
\centering
\begin{tikzpicture}
  % \bl{ponto p}{nivel l}{comprimento}{nome}
  \mesa
  \bl{0}{0}{2}{c}\bl{3}{0}{1}{a}\bl{5}{0}{1}{b}\bl{3}{1}{3}{d}
  \rot{$S_0$}
\end{tikzpicture}
\caption{Um estado.}
\end{figure}

\subsection*{2. Dois estados com uma ação entre eles}
\begin{figure}[H]
\centering
\begin{tikzpicture}
  \begin{scope}\mesa
    \bl{0}{0}{2}{c}\bl{3}{0}{1}{a}\bl{5}{0}{1}{b}\bl{3}{1}{3}{d}\rot{$S_0$}
  \end{scope}
  \draw[-{Stealth}] (3.3,0.5) -- node[above,font=\scriptsize]{move(d,c,0)} (5.1,0.5);
  \begin{scope}[xshift=5.4cm]\mesa
    \bl{0}{0}{2}{c}\bl{3}{0}{1}{a}\bl{5}{0}{1}{b}\bl{0}{1}{3}{d}\rot{$S_1$}
  \end{scope}
\end{tikzpicture}
\caption{Uma transição.}
\end{figure}

\subsection*{3. Tabela de passos (ação, pré-condições, ADD, DEL)}
\begin{table}[H]
\centering\small
\renewcommand{\arraystretch}{1.25}
\begin{tabularx}{\textwidth}{@{}p{2.3cm} X X X@{}}
\toprule
Ação & Pré-condições (verificadas) & ADD & DEL\\
\midrule
$A_1$: move(d,c,0)\newline $t=0$ &
$\clr(d)$; \dots &
$\at(d,0)$, $\on(d,c)$, \dots &
$\at(d,3)$, $\on(d,a)$, \dots\\
\bottomrule
\end{tabularx}
\caption{Execução manual passo a passo.}
\end{table}

\subsection*{4. Fórmulas}
\begin{align}
&\clr(b,t) \leftrightarrow \forall s,l\,\bigl(\cov(b,s,l,t)\rightarrow
  \esp(b,s,t)\bigr) \tag{A5b}
\end{align}
Vínculo causal no texto: $A_1\xrightarrow{\;\clr(a)\;}A_2$. Ordem: $A_1\prec A_2$.

\subsection*{5. Grafo de ordem parcial}
\begin{figure}[H]
\centering
\begin{tikzpicture}[node distance=1.0cm,
  acao/.style={draw,rounded corners,fill=cora!12,minimum width=4.4cm,
               minimum height=0.75cm,font=\small},
  lig/.style={-{Stealth},thick}]
  \node[acao,fill=black!6] (ini) {Início};
  \node[acao,below=of ini] (a1) {$A_1$: move(d,c,0)};
  \node[acao,below=of a1,fill=black!6] (fim) {Fim};
  \draw[lig] (ini) -- node[right,font=\footnotesize]{\ $\clr(d)$} (a1);
  \draw[lig] (a1) -- node[right,font=\footnotesize]{\ \dots} (fim);
\end{tikzpicture}
\caption{Grafo de ordem parcial.}
\end{figure}

\subsection*{6. Comandos e saída de terminal}
\begin{lstlisting}
$ python3 ../bw2cnf_var.py --cenario 1
Gerado: ... variaveis, ... clausulas (T=4)
\end{lstlisting}

\subsection*{7. Tabela simples}
\begin{center}\small
\begin{tabular}{@{}lll@{}}
\toprule
Coluna 1 & Coluna 2 & Coluna 3\\
\midrule
a & b & c\\
\bottomrule
\end{tabular}
\end{center}


\begin{thebibliography}{9}
\bibitem{rn} S. Russell, P. Norvig. \emph{Inteligência Artificial: Uma Abordagem
  Moderna}. Capítulos de planejamento.
\bibitem{bratko} I. Bratko. \emph{Prolog Programming for Artificial Intelligence}.
  Capítulo de planejamento.
\bibitem{manual} E. Mota. \emph{Manual de Resolução do Trabalho: Mundo dos Blocos de
  Tamanho Variável via SAT Solver}. IComp/UFAM, 2026.
% \bibitem{chave} Autor. \emph{Título}. Ano.
\end{thebibliography}
\end{document}
